\section{The agent: key idea and design}
\label{sec:agent}

\paragraph{Key idea.}
Wrongway is a race in which a single well-placed wall can add eight or more steps to a path.
A move only looks good or bad once the opponent's best wall replies have been seen.
We therefore treat the game as a classical adversarial search problem
\citep{knuth1975alphabeta}. The design has two parts:
\begin{enumerate}
  \item \textbf{Search as deep as possible within the time limit.} Every component is judged
        by how much depth it buys per second.
  \item \textbf{A cheap but well-calibrated evaluation} at the leaves. Its parameters are
        learned from self-play rather than hand-tuned.
\end{enumerate}
\Cref{fig:pipeline} shows what happens in one call to \code{choose\_action}.

\begin{figure}[!htbp]
\centering
\begin{tikzpicture}[node distance=3mm and 4mm, every node/.style={font=\small},
  box/.style={draw, rounded corners, align=center, minimum height=9mm, fill=gray!8},
  arr/.style={-{Stealth[length=2mm]}}]
  \node[box] (obs) {Observation\\(positions, walls)};
  \node[box, right=of obs] (bb) {Bitboards\\4 move masks};
  \node[box, right=of bb] (cand) {Root candidates\\moves + 8 walls\\+ defensive walls};
  \node[box, right=of cand] (id) {Iterative deepening\\negamax $\alpha$-$\beta$\\depth $1..10$};
  \node[box, right=of id] (out) {Best move of the\\last finished depth};
  \node[box, below=6mm of id] (leaf) {Leaf: table evaluation\\(learned by TD-leaf)};
  \draw[arr] (obs) -- (bb); \draw[arr] (bb) -- (cand); \draw[arr] (cand) -- (id);
  \draw[arr] (id) -- (out); \draw[arr] (id) -- (leaf);
  \draw[arr] (leaf.east) -| node[pos=0.25, above, font=\scriptsize]{stop at \SI{0.55}{\second}} (out.south);
\end{tikzpicture}
\caption{Decision pipeline for one move. All state is rebuilt from the observation, so the agent
keeps working if the arena restarts it.}
\label{fig:pipeline}
\end{figure}

\subsection{Board representation with bitboards}
\label{sec:bitboard}
Squares are numbered $i = 9r + c$. A set of squares is one Python integer whose bit~$i$ is
square~$i$, as in chess engines such as Stockfish \citep{stockfish}. The walls of a position are
stored as four ``can move'' masks (up, down, left, right): bit~$i$ is set if a pawn on $i$ can
step in that direction. Adding a wall clears two bits in two masks. One step of breadth-first
search (BFS) for all squares at once is then
\begin{equation}
  R' = R \;\lor\; ((R \land U) \gg 9) \;\lor\; ((R \land D) \ll 9)
       \;\lor\; ((R \land L) \gg 1) \;\lor\; ((R \land Ri) \ll 1),
  \label{eq:flood}
\end{equation}
where $R$ is the reached set and $U, D, L, Ri$ are the masks.
Profiling had shown that distance computation took 70\% of the run time. Three further tricks
reduce the number of BFS calls:
\begin{enumerate}
  \item \textbf{Backward layers.} One backward flood from the goal row gives the distance of
        every square (\code{goal\_layers}). It is cached per (goal, wall set) because many tree
        nodes share a wall set.
  \item \textbf{Incremental update after a wall} (\code{dist\_after\_wall}). A wall only
        changes distances ``behind'' its four surrounding squares. We restart the forward flood
        at the step where the path first touches those squares and stop at the first backward
        layer that is unaffected. Re-evaluating a wall therefore costs a few iterations instead
        of a full BFS.
  \item \textbf{Skipping our own BFS.} Walls that do not touch our own shortest path cannot
        lengthen it, so no BFS is run for us.
\end{enumerate}
A unit test checks that the shortcuts match a full BFS on thousands of random positions.

\subsection{Search}
\label{sec:search}
\paragraph{Candidate moves.}
A node offers every legal pawn move. It also offers up to $K = 8$ walls that lie on the
opponent's shortest path and actually lengthen it, ranked by $\Delta d_{\text{opp}} -
\Delta d_{\text{own}}$. Restricting walls to the opponent's path keeps the branching factor near
12 instead of up to 132.
At the root only, two extra kinds of candidate are added:
\begin{itemize}
  \item \emph{Defensive walls.} For each of the three most dangerous walls the opponent could
        place (each costing us at least two steps), we try our own wall that overlaps or crosses
        that position, so the opponent can no longer place it.
  \item \emph{Arena check.} Every root candidate is checked against the arena's
        \code{legal\_moves}/\code{legal\_walls} lists, so an illegal action is impossible.
\end{itemize}

\paragraph{Negamax with iterative deepening.}
The root searches depth $1, 2, \dots, 10$ until a \SI{0.55}{\second} deadline. This leaves margin
for the slower grading CPU and for inter-process communication. The tree checks the clock every
64 nodes and aborts with an exception. If the deadline falls in the middle of an iteration, a
root move that has already been \emph{proved} better than the previous best (which is searched
first) is used rather than discarded.

\paragraph{Move ordering.}
Alpha-beta prunes most when the best move comes first. We use four standard techniques:
\begin{itemize}
  \item a transposition table keyed by (side, both pawns, both wall counts, wall masks) whose
        stored best move is searched \emph{before} the other children are generated. When it
        causes a cutoff, all the wall BFS calls of that node are skipped;
  \item two killer moves per ply \citep{akl1977killer};
  \item the history heuristic \citep{schaeffer1989history} as a tie-breaker;
  \item principal-variation search with null windows \citep{marsland1982pvs}.
\end{itemize}
Together these make the engine $1.3\times$ faster at depth 5--6, with the same move chosen in
45/45 test positions.

\paragraph{Lazy leaf evaluation.}
The ``vulnerability'' term (\cref{sec:eval}) costs about ten BFS calls. At a leaf we first
compute the cheap part. If it is already below $\alpha$, only the opponent's vulnerability could
raise the score, so that term is computed first and the other is skipped when $\alpha$ still
holds (symmetrically for $\beta$).

\paragraph{Root rules found by replay analysis.}
\begin{itemize}
  \item \emph{Progress tie-break.} A bonus of $0.001$ per step gained is added at the root
        only. Without it, positions in which ``every move gets walled later'' produced equal
        scores and the agent shuffled back and forth for 50 turns while the opponent walked home.
  \item \emph{Rush when out of walls.} With no walls left, only forward moves are considered.
        Retreating cannot make the opponent's walls less dangerous.
\end{itemize}

\Cref{alg:root} summarises the root procedure.

\begin{algorithm}[!htbp]
\caption{Root procedure of \code{choose\_action} (simplified)}
\label{alg:root}
\begin{algorithmic}[1]
\State build bitboards and distances from the observation; $T_{\text{end}} \gets \text{now} + 0.55$\,s
\State $C \gets$ moves $\cup$ top-8 walls $\cup$ defensive walls, filtered by the arena's legal lists
\If{a move in $C$ reaches the goal} \Return it \EndIf
\If{no walls left} keep only forward moves in $C$ \EndIf
\State $b \gets C_1$
\For{$d = 1, \dots, 10$}
  \State score each $c \in C$ with $-\textsc{Negamax}(c, d-1)$ $+\ 0.001\cdot\text{progress}(c)$, PVS windows
  \If{deadline hit} use the partially searched winner if it beat $b$; \textbf{break} \EndIf
  \State $b \gets$ best move; move $b$ to the front of $C$
\EndFor
\State \Return $b$ \Comment{any exception $\Rightarrow$ return the shortest-distance legal move}
\end{algorithmic}
\end{algorithm}

\subsection{Evaluation function}
\label{sec:eval}
Let $p$ be the player to move and $q$ the opponent. $d_x$ is the shortest-path distance of
player~$x$ (pawns ignored), $w_x$ their remaining walls, and $u_x$ their
\emph{vulnerability}: the largest number of extra steps that one wall on the first eight steps
of their shortest path would cost them. The score seen by $p$ is
\begin{align}
  E(p) &= \tau + \sigma\, s + V(d_p, d_q, w_p) - V(d_q, d_p, w_q)
          + \nu[w_p]\, u_q - \nu[w_q]\, u_p, \label{eq:eval}\\
  V(d, d', w) &= D[\min(d,40)] + W[w][\min(d',20)] + L[w][\operatorname{clip}(d'-d,-6,6)],
  \label{eq:side}
\end{align}
where $s = +1$ if $w_q = 0$ and $d_p \le d_q$ (an uncatchable race), $s=-1$ in the mirrored
case, and $s = 0$ otherwise.
\Cref{tab:features} lists the tables: 428 parameters in total.
\begin{itemize}
  \item The \emph{lead} table $L$ teaches the agent to keep walls in reserve when it is far
        ahead. An earlier version spent nine walls while leading by 9--15 steps and then lost.
  \item The evaluation is \emph{symmetric}: swapping the players negates $E$ apart from the
        tempo constant $\tau$. Our earlier linear evaluation was not, which produced the
        odd/even bias analysed in \cref{sec:process}.
\end{itemize}

\begin{table}[!htbp]
\centering\small
\caption{Lookup tables of the evaluation (\cref{eq:eval,eq:side}).}
\label{tab:features}
\begin{tabularx}{\linewidth}{@{}l r X@{}}
\toprule
Table & Entries & Meaning \\
\midrule
$D[d]$ & 41 & value of being $d$ steps from the goal \\
$W[w][d']$ & $11\times21$ & value of holding $w$ walls when the opponent is $d'$ steps away \\
$L[w][\ell]$ & $11\times13$ & bonus for a lead of $\ell\in[-6,6]$ steps while holding $w$ walls \\
$\nu[w]$ & 11 & weight of the opponent's vulnerability when we hold $w$ walls \\
$\tau,\ \sigma$ & 2 & tempo for the side to move; safe-race bonus \\
\midrule
Total & 428 & \\
\bottomrule
\end{tabularx}
\end{table}

\subsection{Learning the evaluation}
\label{sec:learning}
\paragraph{Stage 1: evolution of a linear evaluation.}
The first evaluations were linear in seven features. Their weights were found with a $(1+\lambda)$
evolution strategy \citep{beyer2002es}, with $\lambda = 6$ mutants per generation and
self-adaptive step size. A mutant replaced the champion only if it:
\begin{enumerate}
  \item beat the champion in side-swapped games (score at least 0.55);
  \item passed a fresh confirmation match;
  \item was not worse against a pool of older champions and fixed sparring bots.
\end{enumerate}

\paragraph{Stage 2: TD-leaf($\lambda$) on lookup tables.}
The linear formula was rewritten exactly as the tables of \cref{tab:features}, so the starting
point played identically. The tables were then refined with TD-leaf($\lambda$)
\citep{sutton1988td,baxter2000tdleaf}, as follows:
\begin{enumerate}
  \item The agent plays itself using its own negamax at depth 5--6.
  \item Each position is scored as $v = \operatorname{sigmoid}(E/K)$ at the leaf of the
        principal variation.
  \item Every table entry active at that leaf is moved by $\eta\, e_t$, with
        $e_t = \sum_k \lambda^k (v_{t+k+1} - v_{t+k})$, $\lambda = 0.7$.
\end{enumerate}
The run used 20 rounds of 200 games and about 5 minutes per round.
The learned tables can be read directly. For example, the step from distance 2 to 1 is worth
about 2.0 points versus about 1.35 far from the goal, and $\nu[w]$ rises from 0.75 with one wall
to 1.18 with six or seven.
Because the submission may not read files, the tables are written into the
\code{WEIGHTS BEGIN/END} block of the source file by the training script.

\subsection{Robustness}
\begin{itemize}
  \item The whole body of \code{choose\_action} is wrapped in \code{try/except}. Any failure
        returns the legal move with the smallest arena-provided distance, never an exception.
  \item No global state is required.
  \item Peak memory is 4--6\,MB.
  \item The longest move measured through the arena is \SI{0.563}{\second} on our laptop, with
        a \SI{0.55}{\second} budget.
\end{itemize}
